% !TeX program = xelatex
\documentclass[UTF8,zihao=-4]{ctexart}
\usepackage[a4paper,margin=2.4cm]{geometry}
\usepackage{amsmath,booktabs}

\begin{document}

\section*{五、数据分析}

\subsection*{1. 最小偏向角与折射率}

设 $\theta_1$、$\theta_1'$ 分别为汞灯黄光经三棱镜折射后达到最小偏向位置时的左右游标读数，
$\theta_2$、$\theta_2'$ 为取下三棱镜后直射光的左右游标读数。
对每对读数求圆周上的较小夹角 $d(a,b)=\min\{|a-b|,\,360^\circ-|a-b|\}$，
再对两侧所得角差取平均，以减小刻度盘偏心的影响。第 1、2 组按记录顺序配对，
第 3 至第 5 组将直射光的两个游标读数交叉配对。原顺序在这三组给出约 $129^\circ$，
超过顶角为 $60^\circ$ 的棱镜在最小偏向位置的几何上限 $120^\circ$
（此时入射角等于出射角 $i$，且 $\theta_0=2i-\varphi$）。
交叉配对后，三个结果均在 $51^\circ$ 左右，与第 2 组接近。因此采用
\[
  \theta_0=
  \begin{cases}
    [d(\theta_1,\theta_2)+d(\theta_1',\theta_2')]/2, & \text{第 1、2 组},\\
    [d(\theta_1,\theta_2')+d(\theta_1',\theta_2)]/2, & \text{第 3 至第 5 组}.
  \end{cases}
\]
按实验讲义取棱镜顶角 $\varphi=60^\circ$，折射率为
\[
  n=\frac{\sin[(\varphi+\theta_0)/2]}{\sin(\varphi/2)}
   =2\sin\frac{60^\circ+\theta_0}{2}.
\]
以第 2 组为例，左游标角差为 $50^\circ54'$，右游标的读数直接相减为 $309^\circ03'$，
故右游标取较小夹角 $360^\circ-309^\circ03'=50^\circ57'$。
两角平均得到 $\theta_0=50^\circ55.5'$，代入得 $n=1.647511$。

\begin{center}
  \small
  表\,1\quad 游标读数与折射率复算结果\par\smallskip
  \begin{tabular}{@{}ccccccc@{}}
    \toprule
    组次 & $\theta_1$ & $\theta_1'$ & $\theta_2$ & $\theta_2'$ & $\theta_0$ & $n$ \\
    \midrule
    1 & $115^\circ14'$ & $298^\circ20'$ & $125^\circ24'$ & $305^\circ29'$
      & $8^\circ39.5'$ & 1.127893\textsuperscript{*} \\
    2 & $132^\circ47'$ & $312^\circ50'$ & $183^\circ41'$ & $3^\circ47'$
      & $50^\circ55.5'$ & 1.647511 \\
    3 & $146^\circ31'$ & $326^\circ35'$ & $17^\circ37'$ & $197^\circ59'$
      & $51^\circ15'$ & 1.650720 \\
    4 & $151^\circ39'$ & $331^\circ42'$ & $22^\circ39'$ & $202^\circ34'$
      & $50^\circ56'$ & 1.647593 \\
    5 & $164^\circ10'$ & $344^\circ16'$ & $35^\circ15'$ & $215^\circ10'$
      & $50^\circ59.5'$ & 1.648170 \\
    \bottomrule
  \end{tabular}
\end{center}

\noindent\textsuperscript{*} 第 1 组左侧入射光读数的个位有涂改，按 $125^\circ24'$ 转录。
该组左右游标分别给出 $10^\circ10'$ 和 $7^\circ09'$，相差 $3^\circ01'$。
两折射光游标读数相差 $183^\circ06'$，也明显偏离对置游标应有的 $180^\circ$。
这两项异常加上手写偏向角与复算值不一致，表明该组读数需要复测，故仅列入表中备查。
第 3 组交叉配对后，两侧角差为 $51^\circ28'$ 和 $51^\circ02'$，相差 $26'$，
为保留数据中最大，可能与读数或狭缝像对准有关，
其偏差计入重复测量的不确定度。

\subsection*{2. 算术平均值与不确定度}

以第 2 至第 5 组为有效数据，组数 $m=4$。用未舍入的折射率求算术平均值，
并以平均值的实验标准差作为 A 类标准不确定度
\[
  \bar n=\frac{1}{m}\sum_{i=2}^{5}n_i=1.648498\ldots,\qquad
  u_{\mathrm A}(\bar n)
    =\sqrt{\frac{\sum_{i=2}^{5}(n_i-\bar n)^2}{m(m-1)}}
    =7.55\times10^{-4}.
\]

B 类分量考虑棱镜顶角公差和游标读数分辨力。将讲义给出的顶角公差 $\pm5'$ 视为均匀分布的半宽，
得到 $u(\varphi)=5'/\sqrt{3}$。游标读数分辨力为 $1'$，
将每次读数的舍入误差视为独立且在 $\pm0.5'$ 内均匀分布。
每组偏向角由四个读数按绝对值均为 $1/2$ 的系数组合，
故单组的分辨力分量为 $0.5'/\sqrt{3}$，四组平均后为
$u_{\mathrm B}(\bar\theta_0)=0.5'/\sqrt{3m}$。
顶角误差由各组共用，其标准不确定度直接传播到平均折射率。

在 $\varphi=60^\circ$、$\bar\theta_0=51^\circ01.5'$ 处计算灵敏系数。
传播时角度均用弧度，得到
\[
  \begin{aligned}
    c_\theta
      &=\frac{\cos[(\varphi+\bar\theta_0)/2]}{2\sin(\varphi/2)}
        \simeq 0.5662\,\mathrm{rad}^{-1},\\
    c_\varphi
      &=-\frac{\sin(\bar\theta_0/2)}{2\sin^2(\varphi/2)}
        \simeq -0.8614\,\mathrm{rad}^{-1}.
  \end{aligned}
\]
两项传播到折射率的标准不确定度分别约为 $2.38\times10^{-5}$ 和 $7.23\times10^{-4}$。
合成标准不确定度为
\[
  u_{\mathrm c}(\bar n)
    =\sqrt{u_{\mathrm A}^2(\bar n)+
        [c_\theta u_{\mathrm B}(\bar\theta_0)]^2+
        [c_\varphi u(\varphi)]^2}
    \simeq 1.05\times10^{-3}.
\]
因此，三棱镜对汞灯黄光的折射率为
\[
  \boxed{n=1.6485\pm0.0010}\qquad\text{（合成标准不确定度，}k=1\text{）}。
\]

\end{document}
